\section*{Chapitre \ref{ch:rs:measuring-market-making-and-execution} --- Mesurer le market making et l'exécution}

\begin{solution}{exo:rs:measuring-market-making-and-execution:1}
\omterm{def:rs:measuring-market-making-and-execution:pnl}{Capture du spread} $100.00 - 99.99 = +1$ tick~; sélection adverse $99.97 - 100.00 = -3$ ticks~; inventaire $100.02 - 99.97 = +5$ ticks~; le total, $100.02 - 99.99 = +3$ ticks,
est leur somme.
\end{solution}

\begin{solution}{exo:rs:measuring-market-making-and-execution:2}
$8\,144/18\,000 = 45~\%$. Les cotations sont exécutées le plus facilement lorsqu'elles sont mal placées (le prix s'apprête à les franchir), si bien qu'un
\omterm{def:rs:measuring-market-making-and-execution:fill}{taux d'exécution} plus élevé peut signifier davantage de sélection adverse~; le teneur de marché filtré atteint 20~\% et perd moins.
\end{solution}

\begin{solution}{exo:rs:measuring-market-making-and-execution:3}
$0.43 - 1.10 = -0.67$ tick par action, la valeur de la courbe à vingt secondes.
\end{solution}

\begin{solution}{exo:rs:measuring-market-making-and-execution:4}
Pour chaque exécution, les trois termes se télescopent en $m(T) - p$, quelle que soit la valeur de $H$. $H$ déplace la frontière entre sélection adverse et inventaire~: trop court, et le
mouvement défavorable qui se poursuit après $H$ est comptabilisé comme inventaire (la sélection adverse paraît plus faible, l'inventaire pire)~; trop long, et le risque
d'inventaire choisi est comptabilisé comme sélection adverse.
\end{solution}

\begin{solution}{exo:rs:measuring-market-making-and-execution:5}
Délai $80 \times 0.20 = 16$~; exécution $80 \times (50.40 - 50.20) = 16$~; opportunité $20 \times (51.00 - 50.00) = 20$~; frais $0.80$~; \omterm{def:rs:simulation-versus-live:calibration}{implementation shortfall} de 52.80~\$,
soit le gain de 100~\$ du portefeuille théorique moins les 47.20~\$ du portefeuille réel.
\end{solution}

\begin{solution}{exo:rs:measuring-market-making-and-execution:6}
La perte a diminué de 4\,240~\$. Au coût par lot du teneur de marché au meilleur cours (0.85~\$), les 3\,309 lots du teneur filtré auraient perdu 2\,814~\$~: moins d'exécutions
expliquent 4\,111~\$ (97~\%), de meilleures exécutions 129~\$.
\end{solution}

\begin{solution}{exo:rs:measuring-market-making-and-execution:7}
$6.5 \times 252 = 1\,638$ heures, soit 136.5 fois les douze heures simulées~: \omterm{def:rs:measuring-market-making-and-execution:pnl}{capture du spread} $+\$480\,000$, sélection adverse $-1.23$ million de dollars, inventaire $-\$145\,000$, frais
$-\$56\,000$, total $-\$945\,000$. La mise à l'échelle suppose que chaque heure de l'année ressemble aux heures simulées~: même activité, ni ouverture ni clôture, ni nouvelles, ni changement de
concurrents.
\end{solution}

\begin{solution}{exo:rs:measuring-market-making-and-execution:8}
Le shortfall d'un ordre par rapport au \omterm{def:rs:simulation-versus-live:calibration}{prix d'arrivée} est dominé par le mouvement du marché pendant sa durée de vie~: l'ordre de la graine 61 du chapitre a battu le \omterm{def:rs:simulation-versus-live:calibration}{prix
d'arrivée} de 8.69 ticks parce que le cours a baissé, et la moyenne des douze ordres, $-2.00$, a une erreur type de 2.73. Mesurer le prix payé par rapport au mid à chaque
exécution, moyenner sur de nombreux ordres avec leurs erreurs types, conditionner au mouvement du marché et comparer les algorithmes par expérience randomisée.
\end{solution}

\begin{solution}{pb:rs:measuring-market-making-and-execution:1}
\begin{enumerate}
\item 0.43 tick par rapport au mid, 0.19 par rapport au \omterm{def:rs:order-book-features:micro}{microprix}~: le \omterm{def:rs:order-book-features:micro}{microprix} penche vers le mouvement suivant, si bien qu'une partie de la demi-fourchette apparente est
déjà perdue à l'exécution.
\item Informées~: $-1.85$ et $-2.71$. Non informées~: $+0.18$ et $+0.52$.
\item 42~\%.
\item Vingt secondes~: à partir de là, la courbe reste à moins de deux erreurs types de sa valeur à cinq minutes.
\item Cotations du teneur au meilleur cours~: 18\,000 affichées, 8\,144 exécutées, 45~\%~; filtré~: 16\,391, 3\,309, 20~\%.
\item Il augmente la \omterm{def:rs:measuring-market-making-and-execution:pnl}{capture du spread} (0.52 par rapport au mid, 0.44 par rapport au \omterm{def:rs:order-book-features:micro}{microprix}) et le mark-out à vingt secondes ($-0.53$ contre $-0.67$)~;
il ne réduit pas la part des contreparties informées (40~\% contre 42~\%).
\item $-\$0.85$ et $-\$0.81$.
\item \omterm{def:rs:measuring-market-making-and-execution:pnl}{Capture du spread} 3\,519~\$, sélection adverse $-\$8\,977$, inventaire $-\$1\,060$, frais $-\$407$, total $-\$6\,925$.
\item 1\,734~\$, $-\$3\,483$, $-\$771$, $-\$165$, total $-\$2\,685$.
\item Par action, ce sont la valeur de la courbe à l'exécution (0.43) et sa baisse jusqu'à l'horizon ($-1.10$)~; leur somme est le mark-out à vingt secondes.
\item \texteuro1.55 par transaction sur le spread net des frais, une perte de positionnement de \texteuro0.68 (un gain sur les positions de moins de cinq secondes, une perte sur
les plus longues), soit \texteuro0.88 brut.
\item 28\,700 actions exécutées~; délai 0.00, exécution $-7.83$, opportunité $-0.91$, frais 0.05, shortfall $-8.69$ ticks par action~; \omterm{def:rs:measuring-market-making-and-execution:costs}{slippage par rapport au VWAP} $-1.10$.
\item Shortfall $-2.00$ (2.73)~; délai $+2.83$ (0.94)~; exécution $-4.68$ (2.09), dont dérive $-4.78$ (2.07) et prix payé $+0.10$ (0.03)~; opportunité $-0.20$ (0.20)~;
frais 0.05.
\item Le prix payé par rapport au mid à chaque exécution, $+0.10$ tick par action~; le reste relève du marché.
\item $+0.25$ tick (0.24)~: le mid a bougé de $-9.71$ avec l'ordre et de $-9.96$ sans lui. Le prix efficient simulé ignore l'ordre, si bien qu'aucun impact ne dure.
\item \textbf{Résultat énoncé.} Ramené à l'année~: \omterm{def:rs:measuring-market-making-and-execution:pnl}{capture du spread} $+\$480\,000$, sélection adverse $-1.23$ million de dollars, inventaire $-\$145\,000$, frais $-\$56\,000$, total
$-\$945\,000$~; les mark-outs se stabilisent à vingt secondes.
\item Rencontrer moins de traders informés (un signal de fair value qui anticipe le prix efficient, et non la seule file d'attente) ou coter plus large lorsqu'ils sont actifs~;
l'inventaire n'est pas le problème.
\item Par expérience randomisée sur le flux de la société, sur le prix payé par rapport au mid et sur les mark-outs, avec erreurs types (chapitre~21).
\item L'indice de référence, la décomposition, le nombre d'ordres, les erreurs types et ce sur quoi l'on a conditionné.
\item Ce qu'il advient du prix après chaque exécution, moyenné sur les exécutions et les horizons~: quelle part du spread les exécutions conservent.
\end{enumerate}
\end{solution}

\begin{solution}{iq:rs:measuring-market-making-and-execution:1}
Le P\&L par action de l'exécution par rapport au prix de référence à un horizon ultérieur, $s(m(t+h) - p)$~; utiliser une courbe d'horizons, de l'exécution jusqu'au point
où elle se stabilise (vingt secondes pour le teneur du chapitre), par rapport au mid et au \omterm{def:rs:order-book-features:micro}{microprix}.
\end{solution}

\begin{solution}{iq:rs:measuring-market-making-and-execution:2}
\omterm{def:rs:measuring-market-making-and-execution:pnl}{Capture du spread} (la référence à l'exécution moins le prix), sélection adverse (le mouvement de la référence jusqu'à un horizon $H$), inventaire (après $H$), frais~; exact
par construction. Choisir $H$ là où la \omterm{def:rs:measuring-market-making-and-execution:curve}{courbe de mark-out} se stabilise, et rendre compte de la sensibilité à ce choix.
\end{solution}

\begin{solution}{iq:rs:measuring-market-making-and-execution:3}
Il gagne le spread et en rend moins de la moitié sur les positions, soit un profit de \texteuro0.88 par transaction, plusieurs fois par jour~; Menkveld a constaté que les
positions étaient rentables sous cinq secondes et coûteuses au-delà, si bien que son avantage tient à la rapidité de dénouement et que son principal risque est de rester avec de l'inventaire.
\end{solution}

\begin{solution}{iq:rs:measuring-market-making-and-execution:4}
Pas nécessairement~: le \omterm{def:rs:measuring-market-making-and-execution:fill}{taux de touche} augmente lorsque les cotations deviennent plus généreuses que celles de la concurrence, c'est-à-dire lorsque les clients sont les
mieux informés sur le prix. Regarder les mark-outs des transactions remportées.
\end{solution}

\begin{solution}{iq:rs:measuring-market-making-and-execution:5}
Randomiser les ordres (ou les jours-actions) entre eux, mesurer les coûts ancrés à chaque exécution (prix payé par rapport au mid, mark-outs) et le shortfall conditionné au
mouvement du marché, et utiliser des ajustements de type \omterm{def:rs:live-experiments:cuped}{CUPED} ainsi qu'un nombre d'ordres suffisant pour l'erreur type (chapitre~21).
\end{solution}

\begin{solution}{iq:rs:measuring-market-making-and-execution:6}
Entrées~: chaque exécution avec le sens, le prix, la quantité, les horodatages (de la bourse et locaux), les identifiants d'ordre et de stratégie, la contrepartie ou le type de flux lorsqu'il est connu~; les données de marché permettant
de reconstruire le mid et le \omterm{def:rs:order-book-features:micro}{microprix}. Horizons fixes~; regroupements par stratégie, plateforme, contrepartie, heure de la journée, taille. Contrôles~: le P\&L décomposé égal au P\&L comptable
au centime près~; le nombre d'exécutions égal à celui de la bourse~; références prises strictement avant chaque exécution.
\end{solution}
